\documentclass{ctexart}
\usepackage{template}
\usepackage{extarrows}

\begin{document}\pagestyle{empty}
\begin{center}\Large
    北京大学数学科学学院2024-25高等数学A1期中考试
\end{center}
\begin{problem}[1.(16\songti{分})]
    解答下列各题.
    \begin{enumerate}[label=\textbf{(\arabic*)}]
        \item \textbf{(8\songti{分})}\ 若$$\lim_{x\to1}\dfrac{x^2+ax+b}{\sin\left(x^2-1\right)}=\dfrac{4}{9}$$求参数$a,b$的值.
        \item \textbf{(8\songti{分})}\ 设函数$f(x)$在开区间$(c,d)$上连续.试证明:
            对于任意$x_1,x_2,\cdots,x_n\in(c,d)$,存在$\xi\in(c,d)$使得$\displaystyle f(\xi)=\dfrac{1}{n}\sum_{i=1}^{n}f(x_i)$.
    \end{enumerate}
\end{problem}
\begin{solution}
\begin{enumerate}[label=\textbf{(\arabic*)}]
        \item 分母趋于$0$,有限极限要求$1+a+b=0$.此时
            \[x^2+ax+b=(x-1)(x+a+1),\qquad
            \lim_{x\to1}\dfrac{x^2+ax+b}{\sin(x^2-1)}=\dfrac{a+2}{2}=\dfrac49.\]
            解得$a=-10/9,b=1/9$.
        \item 在有限个数$f(x_1),\ldots,f(x_n)$中取最小值$f(x_p)$与最大值$f(x_q)$,则
            \[f(x_p)\leqslant\dfrac1n\sum_{i=1}^nf(x_i)\leqslant f(x_q).\]
            若$x_p=x_q$,直接取$\xi=x_p$;否则由$f$在两点之间的闭区间上连续及介值定理,可取两点之间的$\xi$满足要求.该闭区间包含于$(c,d)$.
    \end{enumerate}
\end{solution}
\begin{problem}[2.(16\songti{分})]
    解答下列各题.
    \begin{enumerate}[label=\textbf{(\arabic*)}]
        \item \textbf{(8\songti{分})}\ 设函数$$f(x)=\sqrt{x^2+1}\arctan x-\ln\left(x+\sqrt{x^2+1}\right)$$求$\di f(x)$.
        \item \textbf{(8\songti{分})}\ 求函数$$y=\dfrac{1}{4\sqrt{2}}\ln\dfrac{x^2+\sqrt{2}x+1}{x^2-\sqrt{2}x+1}-\dfrac{1}{2\sqrt{2}}\arctan\dfrac{\sqrt{2}x}{x^2-1}$$的一阶导数$y'$.
    \end{enumerate}
\end{problem}
\begin{solution}
\begin{enumerate}[label=\textbf{(\arabic*)}]
        \item 直接求导并约去相同项可得
            \[f'(x)=\dfrac{x\arctan x}{\sqrt{1+x^2}}+\dfrac1{\sqrt{1+x^2}}-\dfrac1{\sqrt{1+x^2}}
            =\dfrac{x\arctan x}{\sqrt{1+x^2}},\]
            所以$\di f(x)=\dfrac{x\arctan x}{\sqrt{1+x^2}}\dx$.
        \item 在原表达式有定义的$x\neq\pm1$处,
            \[\begin{aligned}
              y'&=\dfrac1{4\sqrt2}\left(\dfrac{2x+\sqrt2}{x^2+\sqrt2x+1}-\dfrac{2x-\sqrt2}{x^2-\sqrt2x+1}\right)\\
                 &\quad-\dfrac1{2\sqrt2}\dfrac{-\sqrt2(x^2+1)}{(x^2-1)^2+2x^2}\\
                &=\dfrac{1-x^2}{2(x^4+1)}+\dfrac{1+x^2}{2(x^4+1)}=\dfrac1{x^4+1}.
            \end{aligned}\]
    \end{enumerate}
\end{solution}
\begin{problem}[3.(18\songti{分})]
    设函数$$f(x)=\left(x^2-3x+2\right)^{100}\cos\dfrac{\pi x^2}{4}$$
    \begin{enumerate}[label=\textbf{(\arabic*)}]
        \item \textbf{(5\songti{分})}\ 设函数$u(x),v(x)$任意阶可导.对于正整数$n$,写出函数$y=u(x)v(x)$的$n$阶导数的Leibniz公式.
        \item \textbf{(10\songti{分})}\ 对于正整数$n$满足$1\leqslant n\leqslant100$,求$f^{(n)}(1)$.
        \item \textbf{(3\songti{分})}\ 求$f^{(101)}(2)$.
    \end{enumerate}
\end{problem}
\begin{solution}
    \begin{enumerate}[label=\textbf{(\arabic*)}]
        \item $$\left(u(x)v(x)\right)^{(n)}=\sum_{i=0}^{n}C_n^iu^{(i)}(x)v^{(n-i)}(x)$$
        \item 设$u(x)=\left(x^2-3x+2\right)^{100},v(x)=\cos\dfrac{\pi}{4}x^2$,于是$f(x)=u(x)v(x)$.\\
            设$\alpha(x)=(x-1)^{100},\beta(x)=(x-2)^{100}$,于是$u(x)=\alpha(x)\beta(x)$.\\
            注意到$n<100$时,$\alpha^{(n)}(1)=\dfrac{100!}{(100-n)!}(1-1)^{100-n}=0$,于是
            $n<100$时$$u^{(n)}(1)=\sum_{i=0}^{n}C_n^i\alpha^{(i)}(1)\beta^{(n-i)}(1)=0$$
            而$n=100$时$$u^{(100)}(1)=\alpha^{(100)}(1)\beta(1)=100!\cdot(1-2)^{100}=100!$$
            于是$$f^{(n)}(1)=\left\{\begin{array}{l}
                0,1\leqslant n<100\\
                \dfrac{\sqrt{2}}{2}\cdot100!,n=100
            \end{array}\right.$$
        \item 注意到当且仅当$n=100$时$\beta^{(100)}(2)=100!$,否则$\beta^{(n)}(2)=0$.于是
            $$u^{(101)}(2)=101\alpha^{(1)}(2)\beta^{(100)}(2)=101\cdot100(2-1)^{99}\cdot100!=100\cdot101!$$
            又$v^{(1)}(x)=-\dfrac{\pi x}{2}\sin\dfrac{\pi}{4}x^2$,于是$v^{(1)}(2)=0$.\\
            综上可知$$\begin{aligned}
                f^{(101)}(2)
                &= u^{(101)}(2)v(2)+101u^{(100)}(2)v^{(1)}(2) \\
                &= u^{(101)}(2)v(2) \\
                &= -100\cdot101!
            \end{aligned}$$
    \end{enumerate}
\end{solution}
\begin{problem}[4.(16\songti{分})]
    计算下列积分.
    \begin{enumerate}[label=\textbf{(\arabic*)}]
        \item \textbf{(5\songti{分})}\ $\displaystyle A=\int_{0}^{2\pi}\left|\sin x-\cos x\right|\dx$.
        \item \textbf{(3\songti{分})}\ $\displaystyle B=\int_{0}^{2\pi}\sqrt{1+\sin2x}\dx$.
        \item \textbf{(5\songti{分})}\ $\displaystyle I=\int\sqrt{\e^x-1}\dx$.
        \item \textbf{(3\songti{分})}\ $\displaystyle J=\int\dfrac{x\e^x}{\sqrt{\e^x-1}}\dx$.
    \end{enumerate}
\end{problem}
\begin{solution}
    \begin{enumerate}[label=\textbf{(\arabic*)}]
        \item $$\begin{aligned}
                A
                &= \int_{0}^{2\pi}\left|\sqrt{2}\sin\left(x-\dfrac{\pi}{4}\right)\right|\dx \\
                &= \sqrt{2}\int_{-\frac{\pi}{4}}^{\frac{7\pi}{4}}\left|\sin u\right|\di u \\
                &= \sqrt{2}\int_{0}^{2\pi}\left|\sin u\right|\di u\\
                &= 2\sqrt{2}\int_{0}^{\pi}\sin u\di u \\
                &= 2\sqrt{2}\left.\left(-\cos u\right)\right|_{0}^{\pi} \\
                &= 4\sqrt{2}
            \end{aligned}$$
        \item $$\begin{aligned}
                B
                &= \int_{0}^{2\pi}\sqrt{\sin^2x+\cos^2x+2\sin x\cos x}\dx \\
                &= \int_{0}^{2\pi}\left|\sin x+\cos x\right|\dx \\
                &= \sqrt{2}\int_0^{2\pi}\left|\sin\left(x+\dfrac{\pi}{4}\right)\right|\dx \\
                &= 4\sqrt{2}
            \end{aligned}$$
        \item 置$t=\sqrt{\e^x-1}$,则$x=\ln\left(t^2+1\right),\dfrac{\dx}{\di t}=\dfrac{2t}{t^2+1}$.于是
            $$\begin{aligned}
                I
                &= \int\sqrt{\e^x-1}\dx=\int t\cdot\dfrac{2t\di t}{t^2+1} \\
                &= 2\int\dfrac{t^2+1-1}{t^2+1}\di t \\
                &= 2\left(\int\di t-\int\dfrac{\di t}{t^2+1}\right) \\
                &= 2\left(t-\arctan t\right)+C \\
                &= 2\sqrt{\e^x-1}-2\arctan\sqrt{\e^x-1}+C
            \end{aligned}$$
        \item 置$t=\sqrt{\e^x-1}$,则$x=\ln\left(t^2+1\right),\dfrac{\dx}{\di t}=\dfrac{2t}{t^2+1}$.于是
            $$\begin{aligned}
                J
                &= \int\dfrac{x\e^x}{\sqrt{\e^x-1}}\dx =\int\dfrac{\ln\left(t^2+1\right)\cdot\left(t^2+1\right)}{t}\cdot\dfrac{2t\di t}{t^2+1} \\
                &= 2\int\ln\left(t^2+1\right)\di t\\
                &= 2\int x\di t=2\left(xt-\int t\di x\right) \\
                &= 2x\sqrt{\e^x-1}-4\sqrt{\e^x-1}+4\arctan\sqrt{\e^x-1}+C
            \end{aligned}$$
    \end{enumerate}
\end{solution}
\begin{problem}[5.(14\songti{分})]
    解答下列各题.
    \begin{enumerate}[label=\textbf{(\arabic*)}]
        \item \textbf{(5\songti{分})}\ 设函数$f(x)$在$x=a$可导,试证明$$\lim_{h\to0}\dfrac{f(a+h)-f(a-h)}{h}=2f'(a)$$
        \item \textbf{(2\songti{分})}\ 举例说明:即使$f(x)$在$x=a$处连续且$\displaystyle\lim_{h\to0}\dfrac{f(a+h)-f(a-h)}{h}$存在,也不能保证$f'(a)$存在.
        \item \textbf{(5\songti{分})}\ 设函数$f(x)$在$x=a$可导,对于常数$k\neq0,1$,试证明$$\lim_{h\to0}\dfrac{f(a+kh)-f(a-h)}{h}=(k+1)f'(a)$$
        \item \textbf{(2\songti{分})}\ 举例说明:对于常数$k\neq0,1$,即使$\displaystyle\lim_{h\to0}\dfrac{f(a+kh)-f(a-h)}{h}$存在,也不能保证$f'(a)$存在.
    \end{enumerate}
\end{problem}
\begin{solution}
\begin{enumerate}[label=\textbf{(\arabic*)}]
        \item 将分子写成$[f(a+h)-f(a)]+[f(a)-f(a-h)]$,除以$h$并取极限,即得$2f'(a)$.
        \item 取$f(x)=|x-a|$.其在$a$处连续,对称差商恒为$0$,但左右导数分别为$1,-1$,所以$f'(a)$不存在.
        \item 由可导性,
            \[\begin{aligned}
            \dfrac{f(a+kh)-f(a-h)}h
            &=k\dfrac{f(a+kh)-f(a)}{kh}+\dfrac{f(a-h)-f(a)}{-h}\\
            &\longrightarrow(k+1)f'(a).
            \end{aligned}\]
            \textbf{校订说明：}原题右端写为$(k-1)f'(a)$,与题干中的$f(a-h)$不符.按此差商的写法,正确系数是$k+1$;此处已据此修正题干.
        \item 取$f(x)=x-a$（$x\neq a$）,而$f(a)=1$.当$k\neq0$时,
            \[\dfrac{f(a+kh)-f(a-h)}h=k+1,\]
            极限存在,但$f$在$a$处不连续,所以不可导.
    \end{enumerate}
\end{solution}
\begin{problem}[6.(20\songti{分})]
    证明下列各题.
        \begin{enumerate}[label=\textbf{(\arabic*)}]
            \item \textbf{(10\songti{分})}\ 设序列$\left\{x_n\right\}_{n=1}^{\infty}$有极限$\displaystyle\lim_{n\to\infty}x_n=a$.试用序列极限的定义证明$$\displaystyle\lim_{n\to\infty}\dfrac{1}{n}\sum_{i=1}^{n}x_i=a$$.
            \item \textbf{(6\songti{分})}\ 设正数序列$\left\{x_n\right\}_{n=1}^{\infty}$有极限$\displaystyle\lim_{n\to\infty}x_n=a$.试证明$$\lim_{n\to\infty}\sqrt[n]{x_1x_2\cdots x_n}=a$$
            \item \textbf{(4\songti{分})}\ 设序列$\left\{x_n\right\}_{n=1}^{\infty}$满足$\displaystyle\lim_{n\to\infty}\dfrac{1}{n}\sum_{i=1}^{n}x_i=a$,又有$\displaystyle\lim_{n\to\infty}n\left(x_n-x_{n-1}\right)=0$.试证明$\displaystyle\lim_{n\to\infty}x_n=a$.
        \end{enumerate}
\end{problem}
\begin{solution}
\begin{enumerate}[label=\textbf{(\arabic*)}]
        \item 对任意$\ep>0$,取$N_0$,使$i>N_0$时$|x_i-a|<\ep/2$.令$C=\sum_{i=1}^{N_0}|x_i-a|$,则$n>N_0$时
            \[\left|\dfrac1n\sum_{i=1}^nx_i-a\right|\leqslant\dfrac Cn+\dfrac\ep2.\]
            再取$n>2C/\ep$,右端小于$\ep$,故结论成立.
        \item 本小问须补充$x_n>0$（否则乘积的实数$n$次根可能没有定义）.若$a>0$,则$\ln x_n\to\ln a$.由第一问,
            \[\sqrt[n]{x_1\cdots x_n}=\exp\left(\dfrac1n\sum_{i=1}^n\ln x_i\right)\to a.\]
            若$a=0$,对任意$\ep>0$取$N$使$i>N$时$0<x_i<\ep$,记$C=\prod_{i=1}^Nx_i>0$,则
            \[0<\sqrt[n]{x_1\cdots x_n}\leqslant C^{1/n}\ep^{(n-N)/n}\to\ep.\]
            由$\ep$任意可得极限为$0=a$.
        \item 令$S_n=\dfrac1n\sum_{i=1}^nx_i\to a$,并令$t_i=i(x_{i+1}-x_i)$.由题设$t_i\to0$,又有恒等式
            \[x_n-S_n=\dfrac1n\sum_{i=1}^{n-1}i(x_{i+1}-x_i)=\dfrac{n-1}n\cdot\dfrac1{n-1}\sum_{i=1}^{n-1}t_i.\]
            由第一问,右端趋于$0$,所以$x_n=(x_n-S_n)+S_n\to a$.
    \end{enumerate}
\end{solution}
\end{document}
